\subsection{Non-degenerate submodules}

In this no. and the two following, A denotes a ring, S a multiplicative subset of A consisting of elements which are not divisors of zero in A, and B the ring \(S^{-1}A\) ; A is canonically identified with a subring of B ( \(\S\) 2, no. 1, Remark 3). The elements of S are therefore invertible in B.

One of the most important special cases for applications is that where A is an integral domain and S is the set of elements \(\neq0\) of A; B is then the field of fractions of A.

DEFINITION 3. Let M be a sub-A-module of B. M is called non-degenerate if B. M = B.

If B is a field, this condition simply means that M is not reduced to 0.

PROPOSITION 8. Let M be a sub-A-module of B. The following conditions are equivalent:

\begin{enumerate}
\def\labelenumi{(\alph{enumi})}
\item
  M is non-degenerate.
\item
  M meets S.
\item
  If \(j: \mathbf{M} \to \mathbf{B}\) is the canonical injection, the homomorphism \(u = \mathbf{S}^{-1}j: \mathbf{S}^{-1}\mathbf{M} \to \mathbf{B}\) is bijective.
  (a) implies (b), for if B.M = B, there exists \(a \in A\) , \(s \in S\) and \(x \in M\) such that \((a/s)x = 1\) , hence ax = s belongs to S ∩ M. To see that (b) implies (c), note that u is already injective (§2, no. 4, Theorem 1); moreover, if \(x \in M \cap S\) , the image under u of \(x/x \in S^{-1}M\) in B is equal to 1 and u is therefore surjective. Finally, (c) clearly implies (a).
\end{enumerate}

COROLLARY. If M and N are two non-degenerate sub-A-modules of B, the A-modules M + N, M.N and M ∩ N are non-degenerate.

The assertion is trivial for \(\mathbf{M} + \mathbf{N}\) ; on the other hand if \(s \in \mathbf{S} \cap \mathbf{M}\) and \(t \in \mathbf{S} \cap \mathbf{N}\) , then \(st \in \mathbf{S} \cap (\mathbf{M}. \mathbf{N})\) and \(st \in \mathbf{S} \cap (\mathbf{M} \cap \mathbf{N})\) .

Given two sub-A-modules M and N of B, let us denote by N: M the sub-A-module of B consisting of those \(b \in B\) such that \(bM \subset N\) (Chapter I, §2, no. 10, Remark). If every \(b \in N\) : M is mapped to the homomorphism \(h_{b}: x \mapsto bx\) of M to N, a canonical homomorphism \(b \mapsto h_{b}\) is obtained from N: M to \(\operatorname{Hom}_{\mathbf{A}}(\mathbf{M}, \mathbf{N})\) .

PROPOSITION 9. Let M, N be two sub-A-modules of B. If M is non-degenerate, the canonical homomorphism from N: M to \(\mathrm{Hom}_{\mathbf{A}}(\mathbf{M},\mathbf{N})\) is bijective.

Let \(s \in S \cap M\) . If \(b \in N: M\) is such that bx = 0 for all \(x \in M\) , then bs = 0, whence b = 0 since s is not a divisor of 0 in B. On the other hand, let \(f \in \operatorname{Hom}_{\mathbf{A}}(\mathbf{M}, \mathbf{N})\) and set \(b = f(s) / s\) ; for all \(x \in \mathbf{M}\) , there exists \(t \in \mathbf{S}\) such that \(tx \in \mathbf{A}\) . Then

\[
f (x) = s ^ {- 1} t ^ {- 1} f (s t x) = s ^ {- 1} t ^ {- 1} t x f (s) = b x,
\]

whence \(b \in \mathbf{N}\) : M and \(f = h_b\) , which proves the proposition.

Remark. In particular, A: M is canonically identified with the dual M* of M, the canonical bilinear form on M × M* being identified with the restriction to M × (A: M) of the multiplication B × B → B.

